\DocumentMetadata{
  pdfversion=2.0,
  pdfstandard=ua-2,
  lang=en-US,
  tagging=on,
  tagging-setup={math/setup=mathml-SE,tabsorder=structure}
}
\documentclass[11pt,letterpaper]{article}
\usepackage[margin=0.68in,includefoot]{geometry}
\usepackage{newcomputermodern}
\usepackage{amsmath,amscd,mathtools}
\usepackage{tikz,adjustbox}
\providecommand{\square}{\mdlgwhtsquare}
\usepackage[hidelinks]{hyperref}
\hypersetup{
  pdftitle={Analysis PhD Exam, May 2019},
  pdfauthor={University of Florida Department of Mathematics},
  pdfsubject={Graduate mathematics examination},
  pdfdisplaydoctitle=true,
  bookmarksopen=true
}
\setcounter{secnumdepth}{0}
\setlength{\parindent}{0pt}
\setlength{\parskip}{0.28em}
\setlength{\emergencystretch}{3em}
\setlength{\leftmargini}{2.15em}
\setlength{\leftmarginii}{2.65em}
\setlength{\leftmarginiii}{3em}
\pagestyle{empty}
\raggedbottom
\allowdisplaybreaks
\NewTaggingSocketPlug{block/list/label}{examproblem}{%
  \tagstructbegin{tag=Lbl}\tagstructend%
  \tagstructbegin{tag=\LBody}%
  \tagstructbegin{tag=\ProblemHeadingTag,title-o={\ProblemHeadingText}}%
  \tagmcbegin{tag=\ProblemHeadingTag,actualtext={\ProblemHeadingText}}%
  #2%
  \tagmcend%
  \tagstructend%
}
\newenvironment{examproblems}[2]{%
  \def\ProblemHeadingTag{#2}%
  \AssignTaggingSocketPlug{block/list/label}{examproblem}%
  \begin{enumerate}%
  \setcounter{enumi}{#1}%
  \renewcommand{\labelenumi}{\textbf{\theenumi.}}%
  \setlength{\topsep}{0.35em}%
  \setlength{\partopsep}{0pt}%
  \setlength{\itemsep}{0.48em}%
  \setlength{\parsep}{0.18em}%
}{\end{enumerate}}
\newenvironment{examsubparts}{%
  \AssignTaggingSocketPlug{block/list/label}{default}%
  \begin{enumerate}%
  \setlength{\topsep}{0.18em}%
  \setlength{\partopsep}{0pt}%
  \setlength{\itemsep}{0.16em}%
  \setlength{\parsep}{0.1em}%
}{\end{enumerate}}
\newenvironment{examinstructions}{%
  \par\begingroup\itshape
}{%
  \par\endgroup\vspace{0.18em}
}
\makeatletter
\renewcommand\subsection{\@startsection{subsection}{2}{0pt}%
  {-1.25ex plus -.25ex minus -.1ex}%
  {0.55ex plus .1ex}%
  {\normalfont\large\bfseries}}
\renewcommand{\@maketitle}{%
  \newpage\null\vskip -1em%
  {\footnotesize\bfseries
    \href{https://www.ufl.edu/}{University of Florida}, \href{https://math.ufl.edu/}{Department of Mathematics}\par}%
  \vskip -0.5em%
  \tagstructbegin{tag=H1,title={Analysis PhD Exam, May 2019}}%
  \tagpdfparaOff%
  \tagmcbegin{tag=H1}%
    {\LARGE\bfseries\href{https://gma.math.ufl.edu/exams/analysis-phd/}{Analysis PhD Exam}, May 2019\par}%
  \tagmcend%
  \tagpdfparaOn%
  \tagstructend%
  \vskip 0.25em%
  \tagmcbegin{artifact}%
    \hrule height 1pt%
  \tagmcend%
  \vskip 1em%
}
\makeatother
\title{Analysis PhD Exam, May 2019}
\author{}
\date{}
\begin{document}
\maketitle
\thispagestyle{empty}
\begin{examinstructions}
Do six problems.
\end{examinstructions}
\begin{examproblems}{0}{H2}
\def\ProblemHeadingText{Problem 1.}
\item
Let \(X,Y\) be topological spaces. Prove that if \(f:X\to Y\) is continuous, then~\(f\) is Borel measurable.
\def\ProblemHeadingText{Problem 2.}
\item
Rigorously determine the following expressions:
\begin{examsubparts}
\item[\textnormal{(a)}]
\[
\lim_{n\to\infty}\int_0^\infty \frac{\sin(x/n)}{(1+x/n)^n}\,dx
\]
\item[\textnormal{(b)}]
\[
\lim_{n\to\infty}\int_0^\infty \frac{n\sin(x/n)}{x(1+x^2)}\,dx
\]
\item[\textnormal{(c)}]
\[
\lim_{n\to\infty}\int_0^\infty \frac{n}{n^2+x^2}\,dx
\]
\end{examsubparts}
\def\ProblemHeadingText{Problem 3.}
\item
Prove that if a sequence of functions converges in \(L^1\), then it has a subsequence that converges pointwise almost everywhere.
\def\ProblemHeadingText{Problem 4.}
\item
State and prove the Hahn decomposition theorem. For the proof, you may want to consider a sequence of sets such that \(\mu(E_n)\to\sup_E\mu(E)\) rapidly, under the assumption that such a supremum is bounded.
\def\ProblemHeadingText{Problem 5.}
\item
Let~\(X\) be a Banach space. Let \(Y,Z\subseteq X\) be closed subspaces such that \(Y+Z=X\) and \(Y\cap Z=\{0\}\). Prove there is \(C>0\) such that for all \(y\in Y\) and \(z\in Z\),
\[
C(\lVert y\rVert+\lVert z\rVert)\leq \lVert y+z\rVert\leq \lVert y\rVert+\lVert z\rVert.
\]
\def\ProblemHeadingText{Problem 6.}
\item
Let~\(X\) be a Banach space. Show that the natural map from~\(X\) to \(X^{**}\) is injective.
\def\ProblemHeadingText{Problem 7.}
\item
State and prove the Cauchy–Schwarz inequality over a complex Hilbert space.
\def\ProblemHeadingText{Problem 8.}
\item
Prove or give a counterexample: If \(f,g\) are smooth functions on~\(\mathbb{R}\) such that \(f,g\in L^1(\mathbb{R})\), then \(f*g\) is a smooth function and \(f*g\in L^1(\mathbb{R})\).
\end{examproblems}
\end{document}
